\itemhead{4.2 例题 · 例题5（教材p.112）}{XJ-B1-C04-EX-P112-Y0737}
\label{XJ-B1-C04-EX-P112-Y0737}
{\footnotesize 湘教版教材p.112 / PDF第119页 / 4.2 例题 / 例题5}
\field{原题}{已知指数函数 $f(x)=a^x$ 的图象经过点 $(2,7)$，求 $f(-6)$ 和 $f(3)$ 的值。}
\field{参考答案}{$a=\sqrt7$，$f(-6)=1/343$，$f(3)=7\sqrt7$。}
{\small \field{关联知识}{E01 指数函数的概念、图象与性质（教材帮p.150、151 / PDF159、160）；P02 有理数、实数指数幂与运算（教材帮p.144、145 / PDF153、154）}
\field{方法与易错点}{先查底数条件，利用已知点定底数，再据单调性及图象求值或范围。 易错：只要求底数非负；漏排1；将 $a^x+b$ 当标准指数函数且仍用原值域。}
\field{答案依据}{例题原解位置：湘教版教材p.112及题库原图；独立编写可编辑参考解答。答案册此部分主要收练习、习题，不伪造例题答案册位置。}}
